% Generated draft: requires mathematical and source review.
\documentclass[11pt]{article}
\usepackage[margin=25mm]{geometry}
\usepackage{fontspec}
\setmainfont{Latin Modern Roman}
\usepackage{amsmath,amssymb,mathtools,unicode-math}
\setmathfont{Latin Modern Math}
\usepackage{enumitem,needspace,xurl,hyperref}
\setlength{\parindent}{0pt}
\setlength{\parskip}{5pt}
\setlength{\emergencystretch}{3em}
\newcommand{\N}{\mathbb{N}}
\newcommand{\Z}{\mathbb{Z}}
\newcommand{\Q}{\mathbb{Q}}
\newcommand{\R}{\mathbb{R}}
\newcommand{\C}{\mathbb{C}}
\newcommand{\F}{\mathbb{F}}
\newcommand{\A}{\mathbb{A}}
\newcommand{\PP}{\mathbb{P}}
\newcommand{\E}{\mathbb{E}}
\newcommand{\eps}{\varepsilon}
\newcommand{\dd}{\,\mathrm{d}}
\begin{document}
\section*{UnsolvedMath ID: 2100406}
\textbf{Source problem code:} catalogue entry
\subsection{Definitions and mathematical statement}
For a planar billiard table $\Omega\subset\mathbb R^2$ with boundary regular enough for specular reflection, a convex caustic is a convex closed curve inside the table such that a billiard trajectory tangent to it remains tangent after reflection. When tangent trajectories form an invariant circle on which the billiard map is an orientation-preserving circle homeomorphism $f$, its rotation number is represented by
\[
\rho(f)=\lim_{j\to\infty}\frac{\widetilde f^{\,j}(t)-t}{j}\pmod{\mathbb Z},
\]
where $\widetilde f$ is a lift to $\mathbb R$; use the billiard convention taking the relevant value in $[0,1/2]$.

\textbf{Question.} Suppose that a billiard table has a sequence of convex caustics $(C_n)_{n\geq1}$ with rotation numbers $\rho_n$ such that
\[
\lim_{n\to\infty}\rho_n=\omega,\qquad 0<\omega<\tfrac12.
\]
Must its boundary be an ellipse? Here an ellipse includes a circle as the equal-semi-axis case. The source does not state the boundary regularity or strict convexity assumptions, nor whether the caustics or rotation numbers must be distinct. The literal sequence hypothesis permits repetition; an intended accumulating-family version must explicitly exclude that possibility rather than silently strengthening the hypothesis.

The supplied research summary reports ellipse rigidity for a limiting rotation number $\omega=1/2$, but this endpoint is outside the interval in the question. It therefore gives adjacent evidence rather than a resolved subcase of the displayed interior-limit hypothesis. It states that the arbitrary interior-limit case, including irrational limits, has not been verified as settled.
\subsection{Short English statement}
Does a sequence of convex billiard caustics whose rotation numbers approach an interior value force an elliptical table? Distinctness and table regularity are missing, and the cited endpoint progress does not cover the stated interval.
\subsection{Sources}
\begin{itemize}
\item \emph{Open Problems, Questions, and Challenges in Finite-Dimensional Integrable Systems} (2018), Question 4.6, \url{https://arxiv.org/abs/1804.03737}; source-TeX extraction, accessed 2026-07-29. No author list is supplied in this record.
\item The supplied summary reports rigidity at rotation number $1/2$ but supplies no author, publication, theorem statement, or regularity hypotheses for it.
\item Local registry \texttt{data/unsolvedmath/problems.json}, permanent ID $2100406$, source code catalogue entry, \emph{Open Problems in Integrable Systems --- Geometry of caustics, invariant surfaces, and commuting billiard maps}. The record labels this partially solved and PARTIAL-PROGRESS, with NEEDS\_REVIEW source-status and rights qualifications.
\end{itemize}
\end{document}
